\documentclass{uofa-eng-assignment}
\usepackage{caption}
\usepackage{hyperref}
\hypersetup{
    colorlinks = false,
    linkbordercolor = {white},
}
\captionsetup[table]{name=Tabla, labelsep=period}
\captionsetup[figure]{name=Figura, labelsep=period}

\newcommand{\hpl}{\hyperlink}
\newcommand{\hpt}{\hypertarget}
\newcommand*{\name}{Joan Sebastian Ramirez Sanchez  }
\newcommand*{\id}{202223948}
\newcommand*{\course}{Control óptimo aplicado a modelos biológicos}
\newcommand*{\assignment}{Tarea 1 de 3}
\newcommand*{\dated}{\today}
\newcommand{\tf}{\textbf}

\begin{document}

\maketitle
\textbf{Ejercicio 8.2} Resolver
\[
\max_{u} \int_{0}^{1} \bigl(x(t) - u(t)\bigr)\, dt
\]
sujeto a
\[
x'(t) = 2u(t)\bigl(1 - u(t)\bigr), \quad x(0) = 0,
\]
\[
0 \leq u(t) \leq 1.
\]
Como vimos en estas secciones, el problema de este ejercicio radica en la frontera
que tiene nuestro control, es decir, el hecho de que 
$u(t)$ debe estar entre $0$ y $1$. Para resolver este problema, primero encontramos la función Hamiltoniana de nuestro sistema, que es
\[H(x, u, \lambda) = x - u + 2\lambda u(1 - u).\]
Luego, encontramos la función de costo$\lambda(t)$, que es
\[\lambda'(t) = -\frac{\partial H}{\partial x} = -1 \implies \lambda(t) = -t + C.\]
dado que no tenemos condiciones para $x(1)$, podemos hacer que $\lambda(1) = 0$, lo que nos da $C = 1$. Entonces, $\lambda(t) = -t + 1$.
Ahora maximizamos $H$ con respecto a $u$, lo que nos da
\[ H(u) = x-u + 2\lambda u(1 - u) = x - u + 2(-t + 1)u(1 - u) = x - u + 2(-t + 1)(u - u^2).\]
Derivando con respecto a $u$
\[\frac{\partial H}{du} = -1 + 2(-t + 1)(1 - 2u) = -1 + 2(-t + 1) - 4(-t + 1)u = (1-2t) -4u(1-t).\]
Analicemos en los extremos de $u$ para encontrar el máximo.
 \begin{align*}
    \frac{\partial H}{\partial u} <0 &\implies u =0, \implies  (1-2t) < 0, \implies t > \frac{1}{2},
 \end{align*}
esto implica que el máximo del Hamiltoniano se alcanza en $u=0$ cuando $t > \frac{1}{2}$.
\[
\frac{\partial H}{\partial u} >0 \implies u =1, \implies  -3+2t > 0 \implies t > \frac{3}{2},
\]
pero como $t$ solo va hasta $1$, entonces $u = 1$ no es una solución válida. Ahora consideramos el caso interior, cuando
\[
\frac{\partial H}{\partial u} = 0 \implies u = \frac{1-2t}{4(1-t)}.
\]
Este valor es admisible cuando $0 \le u \le 1$, lo cual ocurre para $t \le \frac{1}{2}$.
\begin{align*}
    \begin{cases}
        u = \frac{1-2t}{4(1-t)} \quad  &0 \leq t \leq \frac{1}{2},\\
        u = 0 \quad  &\frac{1}{2} \leq t \leq 1.
    \end{cases}
\end{align*}
\textbf{Ejercicio 8.4}  Lo siguiente es una generalización del Ejemplo 8.4.  
Sean $c, B, T > 0$ constantes y suponga que $B < \frac{cT^2}{4}$. Resolver

\[
\min_{u} \int_0^T \left( u(t)^2 + c\, x(t) \right)\, dt
\]

sujeto a
\[
x'(t) = u(t), \quad x(0) = 0, \quad x(T) = B,
\]
\[
u(t) \geq 0.
\]
Primero, encontramos la función Hamiltoniana de nuestro sistema, que es
\[H(x, u, \lambda) = u^2 + cx + \lambda u.\]
encontramos la función de adjunta
\[\lambda'(t) = -\frac{\partial H}{\partial x} = -c \implies \lambda(t) = -ct + k.\]
para alguna constante $k$. para encontrar $u*$, minimizamos $H$ punto por punto:
\[\frac{\partial H}{\partial u} = 2u + \lambda\]
\[ 
\frac{ \partial H}{\partial u} >0 \implies u = 0, \implies \lambda > 0 \implies -ct + k > 0 \implies t < \frac{k}{c},
\]
\[
\frac{ \partial H}{\partial u} =0 \implies 0 \leq u^* \implies u^* = -\frac{\lambda}{2} = \frac{ct-k}{2}, \quad t \geq \frac{k}{c}.
\]
entonces tenemos la representación de $u^*$:
\[
u^* = \begin{cases}
0 & 0 \leq t < \frac{k}{c},\\
\frac{ct-k}{2} & \frac{k}{c} \leq t \leq T.
\end{cases}
\]
\textbf{Caso 1.} si $k \leq 0$, entonces $0\leq t < \frac{k}{c}$ es imposible,
por lo tanto $x'(t)=u=\frac{ct-k}{2}$, lo que implica que $x(t) = \frac{ct^2}{4} - \frac{kt}{2}$. 
Dado que $x(T) = B=$ implica 
\[
B= \frac{cT^2}{4}-\frac{kT}{2} \implies k = \frac{cT}{2} - \frac{2B}{T}.
\]
\end{document}